\subsection{Coordinación transaccional: 2PC frente a Saga}\label{sec:literatura-coordinacion}
La comparación requiere separar confirmación atómica, aislamiento y recuperación. Garcia-Molina y Salem \cite{garciamolina1987sagas} introducen las sagas como secuencias de transacciones que pueden intercalarse con otras y cuya ejecución parcial se remedia mediante compensaciones. Este fundamento permite distinguir la compensación de una operación ya confirmada del rollback de una transacción todavía abierta. Su objeto original son transacciones de larga duración; no es una evaluación del despliegue de TiendaTech ni una prueba de que cualquier implementación con pasos locales preserve sus invariantes.

Gray y Lamport \cite{gray2006commit} estudian el acuerdo de commit o aborto entre participantes y explican el bloqueo de 2PC ante el fallo de su coordinador. Comparan ese mecanismo con Paxos Commit mediante mensajes, demoras y escrituras persistentes. Para TiendaTech, esto implica que el número de commits locales no basta para representar el coste distribuido: deben observarse comunicación, estado preparado y recuperación. Tampoco debe confundirse el protocolo de confirmación con una garantía de aislamiento independiente del control de concurrencia de sus participantes. La implementación denominada E-2PC carece de esa separación de procesos y no reproduce el modelo completo del artículo.

Štefanko et al.\ \cite{stefanko2019saga} examinan Axon, Eventuate ES, Eventuate Tram y MicroProfile LRA sobre un escenario común, considerando soporte transaccional, complejidad y rendimiento. La evaluación muestra por qué interesa distinguir el patrón de la tecnología concreta que lo ejecuta. Una limitación de su comparación es que el escenario de mayor carga publica la mejor de tres ejecuciones, no una distribución completa. Sus tiempos no son una referencia numérica directamente comparable con el p95 de TiendaTech, ni constituyen un contraste experimental directo entre nuestro E-2PC y E-Saga.

Dürr et al.\ \cite{durr2022saga} amplían la selección de tecnologías de saga con un catálogo de criterios que incluye compensación, persistencia, recuperación, monitorización y pruebas. Su estudio presupone que saga ya resulta adecuada para el caso de uso; no decide si debe preferirse a 2PC. Esta distinción evita utilizar una evaluación de herramientas como argumento para elegir un protocolo. En TiendaTech, una menor latencia debe contrastarse también con la posibilidad de recuperar una compra interrumpida y reconstruir su historial, capacidades que el banco actual no demuestra de extremo a extremo.

Daraghmi et al.\ \cite{daraghmi2022saga} abordan la falta de aislamiento de lectura de saga mediante una propuesta con \emph{quota cache} y un servicio \emph{commit-sync}. El resumen de los autores informa una evaluación sobre comercio electrónico; aquí se utiliza únicamente para identificar el problema y la solución propuesta, sin trasladar cifras ni atribuir garantías generales a su mecanismo. La propuesta modifica el tratamiento del estado antes de persistirlo: no equivale a la saga convencional ni al candado global de nuestro laboratorio. Su pertinencia es recordar que permitir pasos locales no elimina por sí solo las anomalías entre operaciones concurrentes.

En conjunto, estas fuentes delimitan cuatro dimensiones para la comparación: latencia, invariantes observadas, recuperación y coste operativo. La inferencia para este proyecto es que ambos protocolos deben evaluarse bajo los mismos fallos y con participantes independientes, conservando sus controles propios y registrando estados intermedios y finales. No se busca provocar una inconsistencia para confirmar una expectativa: se requiere un instrumento capaz de detectarla si ocurre. La ausencia de anomalías solo es interpretable junto al alcance del oráculo y a las intercalaciones que el ensayo realmente permite.
